% Corrected French for printed p.188; departures are append-only below.
\clearpage
\SourcePage{188}

\LayerTarget{tohoku:remark:ext-bifunctoriality:p1}
\LayerTarget{tohoku:remark:ext-center-modules}
Remarquons enfin, pour finir les sorites, que les
$\mathrm{Ext}_O^p(X;A,B)$ sont des modules sur $\Gamma(O^\natural)$, et
que les $\mathbf{Ext}_O^p(A,B)$ sont des Modules sur le faisceau d'anneaux
$O^\natural$ centre de $O$.

\LayerTarget{tohoku:construction:hom-composite-functors}
La formule (4.1.2) peut aussi s'écrire, en désignant par $h_A$ resp.
$\mathbf{h}_A$ le foncteur covariant de $C^O$ dans la catégorie $C$ des
groupes abéliens, resp. $C^X$, défini par

\LayerTarget{tohoku:eq:4.2.4}
\[
\tag{4.2.4}
h_A(B)=\mathrm{Hom}_O(A,B),\qquad
\mathbf{h}_A(B)=\mathbf{Hom}_O(A,B):
\qquad h_A=\Gamma\mathbf{h}_A.
\]

Donc $h_A$ apparaît comme un foncteur composé, qui est justiciable du
théorème 2.4.1 en vertu de prop. 4.1.3 et du corollaire à la prop. 3.3.1. On
obtient donc :

\LayerTarget{tohoku:theorem:4.2.1}
\noindent\textsc{Théorème 4.2.1.} \textit{Soit $X$ un espace topologique muni
d'un faisceau d'anneaux avec unité $O$, soit $C^O$ la catégorie des
$O$-Modules sur $X$. Soit $A\in C^O$ un $O$-module fixé. Alors il existe sur
$C^O$ un foncteur spectral cohomologique, aboutissant au foncteur gradué
$(\mathrm{Ext}_O^p(X;A,B))$, et dont le terme initial est}

\LayerTarget{tohoku:eq:4.2.5}
\[
\tag{4.2.5}
I_2^{p,q}(A,B)=H^p(X,\mathbf{Ext}_O^q(A,B)).
\]

\LayerTarget{tohoku:eq:4.2.6}
On en conclut en particulier des ``\textit{edge homomorphisms}''
\[
\tag{4.2.6}
\begin{aligned}
H^p(X,\mathbf{Hom}_O(A,B))&\longrightarrow
  \mathrm{Ext}_O^p(X;A,B),\\
\mathrm{Ext}_O^p(X;A,B)&\longrightarrow
  H^0(X,\mathbf{Ext}_O^p(A,B))
\end{aligned}
\]
(ce dernier étant d'ailleurs celui qui résulte de l'interprétation de
$\mathbf{Ext}_O^p(A,B)$ comme le faisceau associé au préfaisceau
$\mathrm{Ext}_O^p(-;A,B)$), et une suite exacte à 5 termes :

\LayerTarget{tohoku:eq:4.2.7}
\[
\tag{4.2.7}
\begin{aligned}
0\longrightarrow H^1(X,\mathbf{Hom}_O(A,B))
&\longrightarrow \mathrm{Ext}_O^1(X;A,B)
\longrightarrow H^0(X,\mathbf{Ext}_O^1(A,B))\\
&\longrightarrow H^2(X,\mathbf{Hom}_O(A,B))
\longrightarrow \mathrm{Ext}_O^2(X;A,B).
\end{aligned}
\]

\LayerTarget{tohoku:remark:ext1-extension-classes}
Cette suite exacte \textit{précise en particulier la structure du groupe
$\mathrm{Ext}_O^1(X;A,B)$ des classes de $O$-Modules extensions de $A$
(Module quotient) par $B$ (sous-Module).}

\LayerTarget{tohoku:construction:stalk-ext-comparison}
Pour pouvoir utiliser la suite spectrale obtenue, il faut préciser encore le
calcul des $\mathbf{Ext}_O^p(A,B)$. Comme le foncteur $F\longmapsto F(x)$ qui,
à un faisceau abélien $F$, associe son groupe ponctuel $F(x)$, est exact, les
$\mathbf{Ext}_O^p(A,B)(x)$ peuvent être considérés, pour $A$ fixé et $p$, $B$
variables, comme formant un foncteur cohomologique covariant universel en $B$,
se réduisant à $\mathbf{Hom}_O(A,B)(x)$ en dimension 0. De même, les
$\mathrm{Ext}_{O(x)}^p(A(x),B(x))$ forment un foncteur cohomologique en $B$,
se réduisant à $\mathrm{Hom}_{O(x)}(A(x),B(x))$ en dimension 0. On conclut
donc, du homomorphisme fonctoriel (4.1.3), (grâce à la définition des
foncteurs cohomologiques universels) des homomorphismes

\LayerTarget{tohoku:eq:4.2.8}
\[
\tag{4.2.8}
\mathbf{Ext}_O^p(A,B)(x)\longrightarrow
\mathrm{Ext}_{O(x)}^p(A(x),B(x))
\]
(qui sont caractérisés par la condition de définir un homomorphisme de
foncteurs cohomologiques, se réduisant à l'homomorphisme (4.1.3) en dimension
0).

\TohokuContinuation{189}
